\documentclass[12pt,a4paper]{article}
\usepackage[margin=25mm, headheight=15pt, headsep=10pt]{geometry}
\usepackage{fontspec}
\usepackage[table]{xcolor} % Note: table option is required for rowcolors
\usepackage{titlesec}
\usepackage{tocloft}
\usepackage{fancyhdr}
\usepackage{booktabs}
\usepackage{tcolorbox}
\tcbuselibrary{skins, breakable}
\usepackage[colorlinks=true, linkcolor=emerald, urlcolor=emerald, bookmarks=true]{hyperref}

\definecolor{emerald}{RGB}{0,102,51}
\definecolor{darkred}{RGB}{139,0,0}
\definecolor{lightgray}{RGB}{245,245,245}

\usepackage[nil]{babel}
\babelprovide[import, main]{bengali}
\babelprovide[import]{arabic}
\babelfont{rm}[Renderer=HarfBuzz,Scale=1.0]{Noto Serif Bengali}
\babelfont{sf}[Renderer=HarfBuzz]{Noto Sans Bengali}
\babelfont{tt}[Renderer=HarfBuzz]{Noto Sans Bengali}
\babelfont[arabic]{rm}[Renderer=HarfBuzz,Scale=1.1]{Noto Naskh Arabic}

% Section styling
\titleformat{\section}{\Large\bfseries\color{emerald}}{\thesection.}{0.5em}{}
\titleformat{\subsection}{\large\bfseries\color{emerald!85!black}}{\thesubsection.}{0.5em}{}
\titleformat{\subsubsection}{\normalsize\bfseries\color{emerald!70!black}}{\thesubsubsection.}{0.5em}{}
\titlespacing*{\section}{0pt}{18pt}{8pt}

% Fancy Header/Footer setup
\pagestyle{fancy}
\fancyhf{} % clear all header and footer fields
\fancyhead[L]{\color{emerald}\small\textbf{\nouppercase{\leftmark}}}
\fancyfoot[L]{\scriptsize\color{black!50}{\lat{AI}}-সংকলিত, আলেম-যাচাইকৃত নয়}
\fancyfoot[C]{\color{emerald!70!black}\thepage}
\renewcommand{\headrulewidth}{0.4pt}
\renewcommand{\headrule}{\hbox to\headwidth{\color{emerald}\leaders\hrule height \headrulewidth\hfill}}
\renewcommand{\footrulewidth}{0pt}

% Custom boxes
% 1. Ayat/Hadith Box
\newtcolorbox{ayatbox}[1][]{
    enhanced, breakable, 
    colback=emerald!5, 
    colframe=emerald, 
    boxrule=0.5pt, 
    arc=3pt, 
    left=10pt, right=10pt, top=10pt, bottom=10pt,
    #1
}

% 2. Quote/Translation Box
\newtcolorbox{quotebox}[1][]{
    enhanced, breakable,
    frame hidden,
    colback=lightgray,
    borderline west={3pt}{0pt}{emerald},
    left=15pt, right=10pt, top=10pt, bottom=10pt,
    sharp corners,
    #1
}

% 3. AI-generated notice box
\newtcolorbox{warnbox}[1][]{
    enhanced, breakable,
    colback=red!4,
    colframe=darkred,
    boxrule=0.8pt,
    arc=3pt,
    left=10pt, right=10pt, top=10pt, bottom=10pt,
    #1
}

% Commands
\newcommand{\arab}[1]{\foreignlanguage{arabic}{#1}}
\newcommand{\kib}[1]{\textcolor{emerald}{\textbf{#1}}}

% Latin (English) fallback: Noto Serif Bengali has NO Latin glyphs
\newfontfamily\latfont[Renderer=HarfBuzz]{Noto Serif}
\newcommand{\lat}[1]{{\latfont #1}}

\setlength{\parskip}{5pt}
\setlength{\parindent}{0pt}

% Fix for missing bullet in Noto Serif Bengali
\renewcommand{\labelitemi}{$\bullet$}



\begin{document}
\begin{center}{\Large\bfseries\color{emerald} আল্লাহর রহমত ও তাওবা — \lat{08}-যুবক-যে-যিনার-অনুমতি-চেয়েছিল}\end{center}
\vspace{1em}
\section*{ঘটনা ৮ — যে যুবক যিনার অনুমতি চেয়েছিল: নবীজির (সা.) স্নেহে পাপ-ইচ্ছার রূপান্তর}

\begin{warnbox}
 \textbf{গুরুত্বপূর্ণ নোটিশ (\lat{Important} \lat{Notice}):} এই কন্টেন্টটি \lat{AI} (কৃত্রিম বুদ্ধিমত্তা) দিয়ে প্রস্তুত ও সংকলিত। \lat{AI} কঠোর সূত্র-মান নীতি মেনে শুধু নির্ভরযোগ্য উৎস থেকে তথ্য সংগ্রহ করেছে — কেবল সহীহ/হাসান হাদিস ও সহীহ সনদের আসার রাখা হয়েছে, দুর্বল সনদ বাদ দেওয়া হয়েছে। তবে এটি কোনো আলেম বা মুহাদ্দিস কর্তৃক যাচাইকৃত নয়।
\end{warnbox}

\begin{quote}
\textbf{সম্পর্কিত ফাইল:} হাদিস ফাইল — তাওবার দরজা; ব্যবহারিক গাইড — শয়তানের কৌশল ও জবাব; কুরআন ফাইল — ২৫:৭০-এর তাফসীর।
\end{quote}

\medskip\hrule\medskip

\subsection*{১. সূত্র কার্ড (\lat{Source} \lat{Card})}

\begin{table}[h]\centering\small
\begin{tabular}{ll}
বিষয় & বিবরণ \\
\hline
\textbf{ঘটনা} & এক যুবক নবীজির (সা.) কাছে এসে যিনার অনুমতি চাইল; উপস্থিতরা ধমক দিতে গেলে নবীজি (সা.) তাকে কাছে ডেকে যুক্তি দিয়ে বুঝালেন, তার জন্য দোয়া করলেন — এরপর যুবকটি আর কোনো হারামের দিকে ফিরে তাকায়নি \\
\textbf{বর্ণনাকারী} & আবু উমামাহ (রাঃ) \\
\textbf{গ্রন্থ ও নম্বর} & মুসনাদ আহমাদ — আবু উমামাহ (রাঃ) থেকে (৫/২৫৬-২৫৭) \\
\textbf{সনদের মান} & \textbf{সহীহ} — আল-ইরাকী (রহ.): সনদ উত্তম, রাবিগণ সহীহের রাবি (\arab{رواه} \arab{أحمد} \arab{بإسناد} \arab{جيد} \arab{رجاله} \arab{رجال} \arab{الصحيح}) — আল-মুগনি আন হামলিল আসফার (২২৬৬) \\
\textbf{স্থান ও সময়} & মদিনা — নবীজির (সা.) জীবদ্দশায় \\
\hline
\end{tabular}
\end{table}

\medskip\hrule\medskip

\subsection*{২. সংক্ষিপ্ত পটভূমি (\lat{Background})}

যুবকের ভেতরে যৌন-ইচ্ছা (\lat{desire}) প্রবল ছিল — ইসলামের বিধান জানার পরও সে যেন নিজেকে সামলাতে পারছিল না। সে নবীজির (সা.) কাছে এসে \textbf{সরাসরি অনুমতি চাইল}। এটি ছিল অত্যন্ত গুরুতর একটি মুহূর্ত — যেকোনো সাধারণ প্রতিক্রিয়ায় এই যুবকটি দ্বীন থেকে দূরে সরে যেতে পারত। নবীজি (সা.) যেভাবে তাকে সামলালেন, তা ইসলামের শিক্ষার (\lat{dawah}) একটি উজ্জ্বল নমুনা — \textbf{পাপীকে ধমক দিয়ে তাড়ানো নয়, বরং কাছে টেনে তার অন্তর পরিবর্তন করা।}

\medskip\hrule\medskip

\subsection*{৩. পূর্ণ ঘটনার বিশুদ্ধ বর্ণনা (\lat{The} \lat{Full} \lat{Narration})}

আবু উমামাহ (রাঃ) বলেন:

\begin{quote}
এক যুবক নবীজির (সা.) কাছে এসে বলল: \textbf{\lat{“}হে আল্লাহর রাসূল! আমাকে যিনার অনুমতি দিন।\lat{”}} উপস্থিত লোকজন তার দিকে ঘুরে ধমক দিতে লাগল — \lat{“}থামো! থামো!\lat{”} নবীজি (সা.) বললেন: \textbf{\lat{“}তাকে কাছে আসতে দাও।\lat{”}} সে কাছে এসে বসল। নবীজি (সা.) জিজ্ঞেস করলেন: \textbf{\lat{“}তুমি কি চাও — এটা তোমার মায়ের সাথে হোক?\lat{”}} সে বলল: \lat{“}না, আল্লাহর কসম! আপনার উপর আমার জীবন উৎসর্গ হোক।\lat{”} নবীজি (সা.) বললেন: \textbf{\lat{“}লোকেরাও তাদের মায়েদের জন্য এটা পছন্দ করে না।\lat{”}} তারপর জিজ্ঞেস করলেন: \textbf{\lat{“}তুমি কি চাও — এটা তোমার মেয়ের সাথে হোক?\lat{”}} সে বলল: \lat{“}না, আল্লাহর কসম!\lat{”} নবীজি (সা.) বললেন: \textbf{\lat{“}লোকেরাও তাদের মেয়েদের জন্য এটা পছন্দ করে না।\lat{”}} এরপর জিজ্ঞেস করলেন: \textbf{\lat{“}তোমার বোনের সাথে?\lat{”}} — \lat{“}না, আল্লাহর কসম!\lat{”} — \textbf{\lat{“}লোকেরাও তাদের বোনেদের জন্য পছন্দ করে না।\lat{”}} তারপর: \textbf{\lat{“}তোমার ফুফুর সাথে?\lat{”}} — \lat{“}না, আল্লাহর কসম!\lat{”} — \textbf{\lat{“}লোকেরাও পছন্দ করে না।\lat{”}} তারপর: \textbf{\lat{“}তোমার খালার সাথে?\lat{”}} — \lat{“}না, আল্লাহর কসম!\lat{”} — \textbf{\lat{“}লোকেরাও পছন্দ করে না।\lat{”}}
\end{quote}

এরপর নবীজি (সা.) তার উপর হাত রাখলেন এবং দোয়া করলেন:

\begin{quote}
\textbf{\arab{اللّٰهُمَّ} \arab{اغْفِرْ} \arab{ذَنْبَهُ} \arab{وَطَهِّرْ} \arab{قَلْبَهُ} \arab{وَحَصِّنْ} \arab{فَرْجَهُ}} \\
\textbf{উচ্চারণ:} আল্লাহুম্মাগফির যামবাহু ওয়া ত্বাহহির ক্বালবাহু ওয়া হাসসিন ফারজাহু \\
\textbf{অর্থ:} \lat{“}হে আল্লাহ! তার পাপ ক্ষমা করো, তার অন্তর পবিত্র করো এবং তার লজ্জাস্থানকে পাপ থেকে রক্ষা করো।\lat{”}
\end{quote}

বর্ণনাকারী বলেন: \textbf{\lat{“}এরপর সেই যুবকটি আর কোনো হারামের দিকে ফিরে তাকায়নি।\lat{”}}

(মুসনাদ আহমাদ — আবু উমামাহ (রাঃ) থেকে; আল-ইরাকীর মতে সনদ সহীহ)

\medskip\hrule\medskip

\subsection*{৪. শব্দের অর্থ (\lat{Key} \lat{Words})}

\begin{table}[h]\centering\small
\begin{tabular}{ll}
শব্দ & অর্থ \\
\hline
\textbf{যিনা (\arab{زنا})} & বিবাহবহির্ভূত যৌন সম্পর্ক — কবিরা গুনাহ \\
\textbf{ইসনাদুন জাইয়িদ (\arab{إسناد} \arab{جيد})} & উত্তম সনদ — সহীহের পর্যায়ে \\
\textbf{ত্বাহহির (\arab{طهّر})} & পবিত্র করো — অন্তর থেকে ময়লা দূর করো \\
\textbf{হাসসিন ফারজাহু (\arab{حصّن} \arab{فرجه})} & তার লজ্জাস্থানকে সুরক্ষিত করো — পাপ থেকে হেফাজত করো \\
\hline
\end{tabular}
\end{table}

\medskip\hrule\medskip

\subsection*{৫. সংশ্লিষ্ট দলিল ও রেফারেন্স}

\begin{itemize}
\item \textbf{দোয়া কবুল ও পরিবর্তন:} দোয়ার পর যুবকটি আর হারামের দিকে তাকায়নি — এটি দেখায় তাওবাকারীর জন্য আল্লাহর সাহায্য বাস্তব; দোয়া ও সংকল্প মিলে মানুষের জীবন বদলে যায়।
\item \textbf{আয়াত (আল-ফুরকান ২৫:৭০):} \lat{“}তারা পাপগুলো নেকিতে বদলে যাবে\lat{”} — এই যুবকের ঘটনা তার বাস্তব নমুনা।
\item \textbf{হাদিস (বুখারি ৭৫০৭; মুসলিম ২৭৫৮):} বারবার পাপ-ক্ষমার হাদিস — আল্লাহর রহমত পাপীর ফেরার অপেক্ষায় থাকে।
\item \textbf{হাদিস (মুসলিম ২৫৭৭):} \lat{“}তোমরা দিনে-রাতে পাপ করো, আর আমি সব পাপ ক্ষমা করি — ক্ষমা চাও।\lat{”}
\end{itemize}
\subsubsection*{মূল শিক্ষা (দলিলভিত্তিক)}

\begin{enumerate}
\item \textbf{পাপীকে ধমক নয় — কাছে টানা:} উপস্থিত সাহাবিরা ধমক দিতে চেয়েছিলেন; নবীজি (সা.) থামিয়ে দিয়ে যুবকটিকে কাছে ডাকলেন। পাপীর মনে সংশোধনের দরজা খোলার জন্য কঠোরতা নয়, বরং স্নেহ (\lat{gentleness}) ও যুক্তি — এটাই নববী দাওয়াহ পদ্ধতি।
\item \textbf{যুক্তি দিয়ে অন্তরে ঢোকা:} নবীজি (সা.) শুধু \lat{“}হারাম\lat{”} বললেন না — তাকে নিজের মা-বোন-মেয়ের কথা ভাবিয়ে ন্যায়সঙ্গতভাবে (\lat{logically}) পাপের কদর্যতা বোঝালেন। একই পাপ, যখন নিজের প্রিয়জনের সাথে কল্পনা করা হয় — তখনই তার ভয়াবহতা স্পষ্ট হয়।
\item \textbf{দোয়া = আল্লাহর হাতে বদল:} নবীজির (সা.) দোয়ার পর যুবকটি সম্পূর্ণ বদলে গেল — মানুষের সংশোধন শেষ পর্যন্ত আল্লাহর তাওফীক (\lat{divine} \lat{guidance}) ছাড়া হয় না; তাই পাপীর জন্য দোয়াই সবচেয়ে বড় সাহায্য।
\item \textbf{এই ঘটনা থেকে একজন পাপীর শিক্ষা:} ইচ্ছা দমনের লড়াই (\lat{struggle}) অস্বাভাবিক নয় — নবীজির (সা.) কাছে আসা, সত্য শোনা, দোয়া চাওয়া — এই তিনটি পদক্ষেপেই জীবন ঘুরে দাঁড়াতে পারে।
\end{enumerate}
\medskip\hrule\medskip

\subsection*{৬. রেফারেন্স}

\begin{itemize}
\item মুসনাদ আহমাদ — আবু উমামাহ (রাঃ) থেকে (৫/২৫৬-২৫৭)
\item আল-মুগনি আন হামলিল আসফার (২২৬৬) — আল-ইরাকীর তাহকীক
\item কুরআন: আল-ফুরকান ২৫:৭০-৭১; আয-যুমার ৩৯:৫৩
\item হাদিস: বুখারি ৭৫০৭; মুসলিম ২৫৭৭, ২৭৫৮
\item ব্যাখ্যা: আবু আমিনা ইলিয়াস — \lat{dailyhadithonline} (\lat{zina} \lat{hadith}); \lat{hadithanswers}.\lat{com} — \lat{source} \lat{verification}
\end{itemize}

\end{document}
